\documentclass[12pt,a4paper]{amsart}

\usepackage{amsmath,amssymb,amsthm}
\usepackage{mathtools}
\usepackage[utf8]{inputenc}
\usepackage[T1]{fontenc}
\usepackage{hyperref}
\usepackage{enumitem,booktabs,array}
\usepackage{geometry}
\geometry{margin=2.5cm}
\usepackage{stmaryrd}

% -------------------------------------------------------
% Theorem environments
% -------------------------------------------------------
\newtheorem{theorem}{Theorem}[section]
\newtheorem{lemma}[theorem]{Lemma}
\newtheorem{corollary}[theorem]{Corollary}
\newtheorem{proposition}[theorem]{Proposition}
\theoremstyle{definition}
\newtheorem{definition}[theorem]{Definition}
\newtheorem{example}[theorem]{Example}
\newtheorem{remark}[theorem]{Remark}
\newtheorem{conjecture}[theorem]{Conjecture}

% -------------------------------------------------------
% Custom commands
% -------------------------------------------------------
\newcommand{\N}{\mathbb{N}}
\newcommand{\Z}{\mathbb{Z}}
\newcommand{\Q}{\mathbb{Q}}
\newcommand{\R}{\mathbb{R}}
\newcommand{\C}{\mathbb{C}}
\newcommand{\F}{\mathbb{F}}
\newcommand{\Zp}{\mathbb{Z}_p}
\newcommand{\Qp}{\mathbb{Q}_p}
\newcommand{\OK}{\mathcal{O}_K}
\newcommand{\p}{\mathfrak{p}}
\newcommand{\Lam}{\Lambda}
\DeclareMathOperator{\Gal}{Gal}
\DeclareMathOperator{\Cl}{Cl}
\DeclareMathOperator{\ord}{ord}
\DeclareMathOperator{\rank}{rank}
\DeclareMathOperator{\Ann}{Ann}
\DeclareMathOperator{\Char}{char}
\newcommand{\cylp}{\mathbb{Q}(\mu_{p^\infty})}

% -------------------------------------------------------
\begin{document}
% -------------------------------------------------------

\title{Iwasawa Theory:\\
       $p$-adic $L$-functions, the Iwasawa Algebra,\\
       and the Main Conjecture}

\author{Michael Fuhrmann}
\address{Specialist Mathematics Project, Build 119}
\email{specialist-maths@localhost}
\date{\today}

\begin{abstract}
Iwasawa theory studies the behaviour of arithmetic objects (class groups,
Selmer groups, Galois cohomology) as one ascends an infinite tower of number
fields.  The central objects are the Iwasawa algebra $\Lambda = \Zp\llbracket T \rrbracket$
(a complete local ring), $p$-adic $L$-functions (constructed analytically
via Kubota--Leopoldt), and Iwasawa modules (finitely generated $\Lambda$-modules
characterised by $\mu$- and $\lambda$-invariants).  We state the Main Conjecture
(proved by Mazur--Wiles 1984 and Wiles 1990) relating the characteristic ideal of the
Iwasawa module to the $p$-adic $L$-function, and sketch the key ingredients
of the proof.  The paper assumes familiarity with algebraic number theory
(rings of integers, class groups, Galois theory).
\end{abstract}

\maketitle
\tableofcontents

% ================================================================
\section{Introduction and Motivation}
% ================================================================

Let $p$ be an odd prime.  Kummer's criterion (1850) states that $p$ divides the
class number $h(\Q(\zeta_p))$ of the $p$-th cyclotomic field if and only if
$p$ is \emph{irregular} (i.e.\ $p$ divides the numerator of a Bernoulli number
$B_{2k}$ for some $k = 1, \ldots, (p-3)/2$).

Iwasawa's insight (around 1959) was to study, not a single cyclotomic field,
but the entire \emph{cyclotomic tower}:
\[
  \Q \subset \Q(\zeta_p) \subset \Q(\zeta_{p^2}) \subset \cdots
  \subset \Q(\zeta_{p^\infty}) = \cylp.
\]
Setting $K_n = \Q(\zeta_{p^n})$, $K_\infty = \bigcup_n K_n$, and
$\Gamma = \Gal(K_\infty/\Q(\zeta_p)) \cong \Zp$ (a pro-$p$ group isomorphic
to the $p$-adic integers), one studies the $p$-part of the class group $A_n$
of $K_n$ as $n$ varies.

\begin{theorem}[Iwasawa growth formula]
\label{thm:growth}
There exist integers $\mu \ge 0$, $\lambda \ge 0$, and $\nu$ (all depending
on $K_\infty/\Q$) such that for all sufficiently large $n$:
\[
  v_p(|A_n|) = \mu p^n + \lambda n + \nu,
\]
where $v_p$ is the $p$-adic valuation.
\end{theorem}

This is the starting point of Iwasawa theory: the growth of $p$-adic class
numbers is governed by just two integers $\mu$ and $\lambda$ (the \emph{Iwasawa invariants}).

% ================================================================
\section{The Iwasawa Algebra}
% ================================================================

\begin{definition}[Iwasawa algebra]
\label{def:iwasawa_algebra}
Let $\Gamma \cong \Zp$ be a topological group (e.g.\ the Galois group of
the cyclotomic $\Zp$-extension).  The \emph{Iwasawa algebra} is the
completed group ring:
\[
  \Lambda = \Zp\llbracket\Gamma\rrbracket
  = \varprojlim_n \Zp[\Gamma/\Gamma^{p^n}]
  \cong \Zp\llbracket T \rrbracket,
\]
the \emph{power series ring} in one variable over $\Zp$, where
$T = \gamma - 1$ for a topological generator $\gamma$ of $\Gamma$.
\end{definition}

\begin{theorem}[Structure of $\Lambda$]
\label{thm:lambda_structure}
The Iwasawa algebra $\Lambda = \Zp\llbracket T \rrbracket$ is:
\begin{enumerate}[label=(\roman*)]
  \item a \emph{complete local ring} with maximal ideal $(p, T)$,
  \item a \emph{regular local ring} of dimension $2$ (Krull dimension),
  \item a \emph{unique factorisation domain},
  \item \emph{Noetherian}.
\end{enumerate}
\end{theorem}

\begin{proof}[Proof sketch]
$\Lambda$ is the $(p, T)$-adic completion of $\Z[T]$, hence complete and
local with residue field $\F_p$.  That it is a UFD follows from the Weierstrass
preparation theorem (see below) and the fact that power series rings over UFDs
are UFDs (Cohen's theorem).
\end{proof}

\begin{theorem}[Weierstrass preparation theorem for $\Lambda$]
\label{thm:weierstrass_prep}
Every non-zero $f \in \Lambda = \Zp\llbracket T \rrbracket$ can be written uniquely as
\[
  f(T) = p^\mu \cdot u(T) \cdot P(T),
\]
where $\mu \ge 0$, $u(T) \in \Lambda^\times$ is a unit, and
$P(T) \in \Zp[T]$ is a \emph{distinguished polynomial} (monic of degree
$\lambda \ge 0$ with all non-leading coefficients divisible by $p$).
The integers $\mu$ and $\lambda$ are the \emph{Iwasawa $\mu$- and
$\lambda$-invariants} of $f$.
\end{theorem}

\begin{remark}
The integers $\mu(f)$ and $\lambda(f) = \deg P$ measure how ``deep''
$f$ is in the maximal ideal.  For the characteristic ideal of an Iwasawa
module, they control the growth in Theorem~\ref{thm:growth}.
\end{remark}

% ================================================================
\section{Finitely Generated $\Lambda$-Modules}
% ================================================================

\begin{theorem}[Structure theorem for $\Lambda$-modules]
\label{thm:structure}
Every finitely generated torsion $\Lambda$-module $M$ is pseudo-isomorphic
to a direct sum:
\[
  M \;\sim\; \bigoplus_{i=1}^s \Lambda/(p^{\mu_i}) \;\oplus\;
             \bigoplus_{j=1}^t \Lambda/(f_j(T)),
\]
where $f_j$ are irreducible distinguished polynomials in $\Zp[T]$.
\emph{Pseudo-isomorphic} means there is a $\Lambda$-module map with finite
kernel and cokernel.
\end{theorem}

\begin{definition}[Characteristic ideal and Iwasawa invariants]
\label{def:char_ideal}
For a finitely generated torsion $\Lambda$-module $M$ pseudo-isomorphic
to the module above, the \emph{characteristic ideal} is
\[
  \mathrm{char}(M) = \Bigl(p^{\mu} \cdot \prod_j f_j(T)\Bigr) \subseteq \Lambda,
\]
where $\mu = \sum_i \mu_i$.
The \emph{$\mu$-invariant} $\mu(M) = \sum_i \mu_i$ and
\emph{$\lambda$-invariant} $\lambda(M) = \sum_j \deg f_j$ are well-defined.
\end{definition}

% ================================================================
\section{The Cyclotomic $\Zp$-Extension}
% ================================================================

\begin{definition}[Cyclotomic $\Zp$-extension]
\label{def:cyclotomic_ext}
Let $F$ be a number field.  The \emph{cyclotomic $\Zp$-extension} of $F$
is $F_\infty = F(\mu_{p^\infty})^{\mathrm{conn}}$, the unique $\Zp$-extension
of $F$ contained in $F(\mu_{p^\infty})$.  For $F = \Q$:
\[
  \Q_\infty = \bigcup_{n \ge 1} \Q(\zeta_{p^n})^+
  \quad\text{(the real subfield of $\Q(\zeta_{p^n})$)}.
\]
The Galois group $\Gamma = \Gal(\Q_\infty/\Q) \cong \Zp$.
\end{definition}

\begin{remark}[The diamond operators]
The Galois group $\Gal(\Q(\zeta_p)/\Q) \cong (\Z/p\Z)^\times$ acts on
the tower via the \emph{Teichmüller character}
$\omega : (\Z/p\Z)^\times \to \mu_{p-1} \subset \Zp^\times$.
The decomposition $\Gal(\Q(\zeta_p)/\Q) = \langle\sigma\rangle \times \Gamma_0$
(with $\Gamma_0 \cong \Zp$ and $|\langle\sigma\rangle| = p-1$) allows one to
decompose Iwasawa modules into \emph{eigenspaces} $e_i M$ for each character
$\omega^i$ ($i = 0, 1, \ldots, p-2$).
\end{remark}

\begin{definition}[The Iwasawa module $X$]
Let $M_n$ be the maximal abelian $p$-extension of $K_n = \Q(\zeta_{p^n})$
unramified outside $p$.  Define
\[
  X = \Gal(M_\infty / K_\infty)
  = \varprojlim_n \Gal(M_n/K_n),
\]
where $M_\infty = \bigcup_n M_n$.  Then $X$ is a finitely generated
torsion $\Lambda$-module (this is a non-trivial theorem).
\end{definition}

% ================================================================
\section{The Kubota--Leopoldt $p$-adic $L$-Function}
% ================================================================

\begin{definition}[Dirichlet $L$-functions at negative integers]
For a Dirichlet character $\chi$ of conductor $f$ with $\chi(-1) = -1$
(odd character), the values $L(1-n, \chi)$ for $n = 1, 2, 3, \ldots$ are
algebraic numbers given by
\[
  L(1-n, \chi) = -\frac{B_{n,\chi}}{n},
  \qquad
  B_{n,\chi} = f^{n-1} \sum_{a=1}^{f} \chi(a) B_n\!\Bigl(\frac{a}{f}\Bigr),
\]
where $B_n(x)$ are Bernoulli polynomials.
\end{definition}

\begin{theorem}[Kubota--Leopoldt, 1964]
\label{thm:kubota_leopoldt}
Let $\chi$ be a primitive Dirichlet character of conductor $f$ with
$\chi(-1) = -1$ (or $\chi = \mathbf{1}$, the trivial character for the
$\zeta$-function case).  Fix an odd prime $p$ not dividing $f$.
There exists a unique continuous function
\[
  L_p(s, \chi) : \Zp \longrightarrow \Qp,
\]
called the \emph{Kubota--Leopoldt $p$-adic $L$-function}, satisfying
\[
  L_p(1-n, \chi) = \bigl(1 - \chi(p)\omega^{-n}(p) p^{n-1}\bigr)
                   L(1-n, \chi\omega^{-n})
\]
for all positive integers $n \geq 1$, where $\omega$ is
the Teichmüller character.
\end{theorem}

\begin{proof}[Construction sketch]
\textbf{Step 1: The $p$-adic measure.}
The key is to construct a \emph{$p$-adic measure} $\mu_\chi$ on $\Zp^\times$
(a finitely additive bounded $\Qp$-valued function on open compact subsets)
such that
\[
  \int_{\Zp^\times} x^{n-1} \, d\mu_\chi(x)
  = (1 - \chi(p)\omega^{-n}(p)p^{n-1}) L(1-n, \chi\omega^{-n}).
\]

\textbf{Step 2: $p$-adic interpolation.}
The values $(1 - \chi(p)p^{n-1}) L(1-n, \chi)$ for $n \equiv 0 \pmod{p-1}$
are $p$-adically close for $n$ in arithmetic progressions.  More precisely,
for $n \to n'$ in $\Zp$: there exists a power series
$G_\chi(T) \in \Zp\llbracket T \rrbracket$ such that
$G_\chi((1+p)^{1-n} - 1) = $ the interpolated values.

\textbf{Step 3: Iwasawa power series.}
Via the identification $\Gamma \cong \Zp$ (sending a topological generator
$\gamma$ of $\Gamma = \Gal(\Q_\infty/\Q)$ to $1+p$), one obtains
\[
  L_p(s, \chi) = G_\chi((1+p)^s - 1) \in \Lambda \otimes_{\Zp} \Qp.
\]
The construction was carried out by Kubota--Leopoldt (1964) using
Stickelberger elements; a cohomological approach was given by Coleman (1979).
\end{proof}

\begin{remark}[Exceptional zero]
For $\chi = \omega^k$ (a power of the Teichmüller character), $L_p(s, \omega^k)$
may vanish at $s = 1-k$ even when $L(1-k, \omega^k) \ne 0$.  This is the
\emph{exceptional zero} (or \emph{trivial zero}) phenomenon, studied by
Ferrero--Greenberg and later by Greenberg--Stevens.
\end{remark}

% ================================================================
\section{The Main Conjecture (Mazur--Wiles Theorem)}
% ================================================================

\begin{conjecture}[Iwasawa Main Conjecture, original form]
\label{conj:main}
Let $p$ be an odd prime.  For each even integer $i$ with $0 \le i \le p-3$,
let $X^{(i)}$ denote the $\omega^i$-eigenspace of $X$, and let
$f_i(T) \in \Lambda$ be a generator of the characteristic ideal
$\mathrm{char}(X^{(i)})$.  Then
\[
  f_i(T) \;\sim\; G_{\omega^{1-i}}(T) \pmod{\mathrm{units}},
\]
where $G_{\omega^{1-i}}(T)$ is the Iwasawa power series of the $p$-adic
$L$-function for the character $\omega^{1-i}$.
\end{conjecture}

\begin{theorem}[Mazur--Wiles 1984; Wiles 1990]
\label{thm:mazur_wiles}
The Iwasawa Main Conjecture is true for all odd primes $p$ and all
even characters $\omega^i$ as above.
\end{theorem}

\begin{proof}[Proof sketch (Wiles' approach)]
Wiles' 1990 proof (which completed and clarified the Mazur--Wiles 1984
argument) proceeds in three major steps:

\textbf{Step 1: Norm-compatible cyclotomic units.}
One constructs norm-compatible elements $c_n \in K_n^\times$ from
\emph{cyclotomic units} (related to circular units and Gauss sums)
whose norms are compatible across the tower.
The key source is $\eta_n = $ the image of $(1 - \zeta_{p^n})$ in $K_n^\times$.

\textbf{Step 2: Index formula for circular units.}
Let $C_n$ be the group of circular units in $K_n$.  One shows
\[
  [E_n : C_n] = h_n^+ \quad\text{(the real class number)},
\]
using the analytic class number formula and the Dirichlet $L$-function values.
This provides the analytic side.

\textbf{Step 3: Hecke algebras, modular curves, and the Fitting ideal comparison.}
Wiles (1990) established the key divisibility
\[
  \mathrm{char}(X^{(i)}) \mid (G_{\omega^{1-i}}(T)) \quad \text{in } \Lambda
\]
using \emph{Hecke algebras and modular curves}, extending the strategy of
Mazur--Wiles (1984): one relates the Iwasawa module $X^{(i)}$ to the
$p$-adic Hecke algebra acting on the Jacobian of the modular curve $X_1(p^n)$,
following Ribet's proof of the converse of Herbrand's theorem.
The reverse divisibility (equality) follows from comparing $\mu$-invariants:
Ferrero--Washington (1979) proved $\mu(X^{(i)}) = 0$ for $\Q$-extensions,
and the $\mu$-invariant of $G_{\omega^{1-i}}$ is also $0$ (by the non-vanishing
of $L(1, \omega^{1-i}) \ne 0$).  With equal $\mu$-invariants, a standard
rank argument forces equality of characteristic ideals.

\textit{Note:} The Euler system approach (using norm-compatible circular units
directly) was developed by Rubin (1991) for imaginary quadratic fields; it
yields a more streamlined proof but was not the method used by Wiles in 1990.
\end{proof}

\begin{remark}[Later generalisations]
\begin{itemize}
  \item Rubin (1991) gave a shorter proof using Euler systems for imaginary
    quadratic fields, establishing the Main Conjecture over $\Q$ and more.
  \item Kato (2004) proved one divisibility of the Main Conjecture for
    elliptic curves over $\Q$ using Beilinson--Kato elements.
  \item The \emph{non-commutative} Main Conjecture for $p$-adic representations
    was formulated by Coates--Fukaya--Kato--Sujatha--Venjakob (2005) and
    proved in special cases.
\end{itemize}
\end{remark}

% ================================================================
\section{The $\mu$- and $\lambda$-Invariants}
% ================================================================

\begin{definition}[$\mu$- and $\lambda$-invariants of a number field]
\label{def:mu_lambda}
For the cyclotomic $\Zp$-extension $K_\infty/K$, the $\mu$- and
$\lambda$-invariants are those of the Iwasawa module $X$ (the Galois group
of the maximal abelian $p$-extension unramified outside $p$):
\[
  \mu(K/p) = \mu(X), \qquad \lambda(K/p) = \lambda(X).
\]
By Theorem~\ref{thm:growth}, $v_p(|A_n|) = \mu p^n + \lambda n + \nu$
for large $n$.
\end{definition}

\begin{conjecture}[Iwasawa $\mu = 0$ conjecture]
For every number field $K$ and every prime $p$, the Iwasawa $\mu$-invariant
of the cyclotomic $\Zp$-extension satisfies $\mu(K_\infty/K) = 0$.
\end{conjecture}

\begin{theorem}[Ferrero--Washington, 1979]
For $K = \Q$ and any odd prime $p$: $\mu(\Q_\infty/\Q) = 0$.
More generally, for any abelian extension $K/\Q$: $\mu(K_\infty/K) = 0$.
\end{theorem}

\begin{proof}[Proof idea]
Ferrero and Washington proved that the $p$-adic $L$-function $L_p(s, \chi)$
does not vanish identically modulo $p$, which forces $\mu = 0$.
Their argument uses the theory of $p$-adic measures and the explicit
form of Bernoulli numbers.
\end{proof}

\begin{remark}
For non-abelian extensions, the $\mu = 0$ conjecture remains open in general.
It is known for CM fields (Gillard, Schneps 1987) and for some GL$_2$-extensions.
\end{remark}

% ================================================================
\section{BSD and Iwasawa Theory for Elliptic Curves}
% ================================================================

\begin{definition}[Selmer group and Iwasawa module for $E/\Q$]
\label{def:selmer_iwasawa}
Let $E/\Q$ be an elliptic curve and $p$ an odd prime of good ordinary
reduction (i.e.\ $a_p(E) \not\equiv 0 \pmod{p}$).  Let $E[p^\infty]$
denote the $p$-power torsion points.  The \emph{$p$-adic Selmer group}
over $\Q_\infty$ is
\[
  \mathrm{Sel}_{p^\infty}(E/\Q_\infty)
  = \ker\Bigl(H^1(\Q_\infty, E[p^\infty]) \to
    \prod_v H^1(\Q_{\infty,v}, E)[p^\infty]/E(\Q_{\infty,v})\Bigr).
\]
The Pontryagin dual $X_E = \mathrm{Sel}_{p^\infty}(E/\Q_\infty)^\vee$
is a finitely generated $\Lambda$-module.
\end{definition}

\begin{conjecture}[Iwasawa Main Conjecture for elliptic curves]
\label{conj:ec_main}
For $E/\Q$ with good ordinary reduction at $p$, there exists a
$p$-adic $L$-function $L_p(E, s) \in \Lambda$ (the Mazur--Swinnerton-Dyer
$p$-adic $L$-function) such that
\[
  \mathrm{char}(X_E) = (L_p(E, s)) \quad \text{in } \Lambda.
\]
\end{conjecture}

\begin{theorem}[Kato 2004; Skinner--Urban 2014]
\label{thm:kato_su}
One divisibility $\mathrm{char}(X_E) \mid (L_p(E, s))$ holds under mild
hypotheses (Kato 2004).  Under further hypotheses (Skinner--Urban 2014
established the other divisibility for a large class of elliptic curves):
the Main Conjecture holds, and in particular,
\[
  L_p(E, 1) = 0 \implies \mathrm{rank}(E(\Q)) \ge 1.
\]
\end{theorem}

\begin{remark}[Connection to BSD]
The Iwasawa Main Conjecture for $E$ implies the $p$-part of the
Birch--Swinnerton-Dyer conjecture: $\ord_{s=1} L(E,s) = \mathrm{rank}(E(\Q))$
and the precise formula for the leading coefficient (up to $p$-adic units).
The full BSD conjecture remains open.
\end{remark}

% ================================================================
\section{Explicit Examples}
% ================================================================

\begin{example}[$p=3$, $\Q(\zeta_3)$]
The prime $p=3$, $\zeta_3 = e^{2\pi i/3}$.
$K_1 = \Q(\zeta_3)$, $K_2 = \Q(\zeta_9)$, $K_3 = \Q(\zeta_{27})$, \ldots
$h_1 = h(\Q(\zeta_3)) = 1$ (class number).
$h_2 = h(\Q(\zeta_9)) = 1$.
$h_3 = h(\Q(\zeta_{27})) = 1$.
By the growth formula, $v_3(h_n) = 0 = \mu \cdot 3^n + \lambda n + \nu$
forces $\mu = \lambda = 0$ (the first non-trivial irregular prime for $p=3$
would be where $3 | h$; none occurs in this tower).
\end{example}

\begin{example}[Irregular prime $p = 37$]
$p = 37$ is the smallest irregular prime: $37 \mid B_{32}$
(the 32nd Bernoulli number).  The class number $h(\Q(\zeta_{37})) = 37$.
The Iwasawa $\lambda$-invariant satisfies $\lambda \ge 1$.
Computations (Buhler--Crandall 1992) give $\mu = 0$, $\lambda = 2$
for this prime.
\end{example}

% ================================================================
\section{Summary and Perspectives}
% ================================================================

\begin{remark}[What Iwasawa theory achieves]
The Main Conjecture (now theorem) establishes a profound dictionary:
\begin{center}
\begin{tabular}{l|l}
\textbf{Analytic side} & \textbf{Algebraic side} \\
\hline
$p$-adic $L$-function $L_p(s, \chi)$ & Characteristic ideal of $X^{(i)}$ \\
$L(1-n, \chi) \ne 0$ & $X^{(i)}$ is finite \\
$\mu(G_\chi) = 0$ (Ferrero--Washington) & $\mu(X^{(i)}) = 0$ \\
$\lambda(G_\chi)$ & rank of $X^{(i)}$ over $\Lambda$ \\
\end{tabular}
\end{center}
\end{remark}

\begin{remark}[Open problems]
\begin{itemize}
  \item \textbf{$\mu = 0$ conjecture} for non-abelian extensions.
  \item \textbf{Non-commutative Main Conjecture} for GL$_n$ motives.
  \item \textbf{$p$-adic BSD}: Full BSD from Iwasawa theory requires
    control of the Tate--Shafarevich group.
  \item \textbf{$p$-adic Langlands correspondence}: The Breuil--Colmez
    programme aims to encode Iwasawa modules via $(\varphi,\Gamma)$-modules.
\end{itemize}
\end{remark}

\section*{Acknowledgements}
The author thanks the Specialist Mathematics Project for providing the
research infrastructure.

\begin{thebibliography}{99}

\bibitem{Iwasawa1959}
K.~Iwasawa,
\textit{On $\Gamma$-extensions of algebraic number fields},
Bull.\ Amer.\ Math.\ Soc.\ \textbf{65} (1959), 183--226.

\bibitem{KubotaLeopoldt1964}
T.~Kubota and H.-W.~Leopoldt,
\textit{Eine $p$-adische Theorie der Zetawerte, I: Einführung der $p$-adischen Dirichletschen $L$-Funktionen},
J.\ Reine Angew.\ Math.\ \textbf{214/215} (1964), 328--339.

\bibitem{MazurWiles1984}
B.~Mazur and A.~Wiles,
\textit{Class fields of abelian extensions of $\Q$},
Invent.\ Math.\ \textbf{76} (1984), 179--330.

\bibitem{Wiles1990}
A.~Wiles,
\textit{The Iwasawa conjecture for totally real fields},
Ann.\ of Math.\ (2) \textbf{131} (1990), 493--540.

\bibitem{Rubin1991}
K.~Rubin,
\textit{The ``main conjecture'' of Iwasawa theory for imaginary quadratic fields},
Invent.\ Math.\ \textbf{103} (1991), 25--68.

\bibitem{FerWash1979}
B.~Ferrero and L.~C.~Washington,
\textit{The Iwasawa invariant $\mu_p$ vanishes for abelian number fields},
Ann.\ of Math.\ (2) \textbf{109} (1979), 377--395.

\bibitem{Kato2004}
K.~Kato,
\textit{$p$-adic Hodge theory and values of zeta functions of modular forms},
Ast\'erisque \textbf{295} (2004), 117--290.

\bibitem{SkinnerUrban2014}
C.~Skinner and E.~Urban,
\textit{The Iwasawa main conjecture for $\mathrm{GL}_2$},
Invent.\ Math.\ \textbf{195} (2014), 1--277.

\bibitem{Neukirch1999}
J.~Neukirch,
\textit{Algebraic Number Theory},
Springer, Berlin, 1999.

\bibitem{Washington1997}
L.~C.~Washington,
\textit{Introduction to Cyclotomic Fields},
2nd ed., Springer, New York, 1997.

\bibitem{Coates2006}
J.~Coates and R.~Sujatha,
\textit{Cyclotomic Fields and Zeta Values},
Springer, Berlin, 2006.

\end{thebibliography}

\end{document}
